% Section: Hyperparameter Optimization for Non-Stationary Routing
% Include from paper appendix

\section{Hyperparameter Optimization for Non-Stationary Routing}
\label{sec:hparam_sweep}

\paragraph{The challenge.}
ParetoBandit's core hyperparameters---exploration coefficient
$\alpha$, prior strength $n_{\text{eff}}$, and forgetting factor
$\gamma$---interact non-trivially with the stationarity of the
deployment environment.
Standard cross-validation assumes i.i.d.\ data and optimises a
single stationary metric (e.g., reward or Pareto AUC), which
selects near-stationary forgetting ($\gamma \geq 0.999$)
because forgetting \emph{always} hurts stationary performance.
However, $\gamma{=}1.0$ fails catastrophically under distribution
shift: as shown in Section~\ref{sec:eval_nonstationary}, the
no-forgetting Naive Bandit accumulates $34\%$ more cumulative regret
than the tuned ParetoBandit ($\gamma{=}0.997$) after a reward swap.
Conversely, optimising solely for adaptability yields
$\gamma \ll 1$ with excessive forgetting that destabilises the
stationary policy.
This creates a fundamental tension absent from supervised routing:
\emph{the same hyperparameter must simultaneously protect stationary
efficiency and enable non-stationary recovery}.

Existing bandit-based LLM routers do not address this tension.
PILOT~\cite{panda2025pilot} grid-searches the exploration rate
$\alpha$ to maximise a single reward metric on held-out data,
with no forgetting mechanism and no non-stationarity consideration
in the selection protocol.
BaRP~\cite{wang2025barp} and LLM Bandit~\cite{li2025llmbandit}
omit forgetting entirely.
MixLLM~\cite{wang2025mixllm} delegates adaptation to periodic
neural-network retraining rather than tuning bandit-specific
parameters.
In the general bandit literature, the CDT framework~\cite{kang2024cdt}
proposes online hyperparameter tuning via a double-layer bandit that
treats $\alpha$ selection as a non-stationary continuum-armed problem,
extending earlier discrete methods~\cite{bouneffouf2020oplinucb}.
However, CDT remains \emph{single-objective} (maximise cumulative
reward) and does not tune the forgetting factor $\gamma$ or address
the efficiency--adaptability tradeoff.
The foundational Discounted UCB analysis by
Garivier and Moulines~\cite{garivier2011switching} derives $\gamma$
theoretically from the number of breakpoints, but this quantity is
unknown in every production deployment.
To the best of our knowledge, no prior work---in LLM routing or in
the broader non-stationary bandit literature---proposes a
multi-objective protocol for jointly selecting $(\alpha, n_{\text{eff}},
\gamma)$ to balance stationary efficiency against non-stationary
recovery.

\paragraph{Adaptation-horizon coupling.}
Rather than treating $\alpha$, $n_{\text{eff}}$, and $\gamma$ as three
independent hyperparameters, we observe that two of
them---$n_{\text{eff}}$ and $\gamma$---are linked through a single
interpretable quantity: the \emph{adaptation horizon} $T_{\text{adapt}}$.
Under the Discounted LinUCB update~\cite{russac2019weighted}, the
effective weight of online evidence relative to the prior reaches
parity after
%
\begin{equation}\label{eq:t_adapt}
  T_{\text{adapt}} = \frac{-\log\bigl(n_{\text{eff}}\,(1{-}\gamma) + 1\bigr)}{\log \gamma}
\end{equation}
%
queries.  Inverting for $n_{\text{eff}}$ gives
$n_{\text{eff}} = (\gamma^{-T_{\text{adapt}}} - 1) / (1 - \gamma)$,
which for $\gamma{=}1.0$ reduces to $n_{\text{eff}} = T_{\text{adapt}}$
by L'H\^opital's rule.
Fixing $T_{\text{adapt}}$ collapses the 3D grid
$(\alpha, n_{\text{eff}}, \gamma)$ to a 2D search over
$(\alpha, \gamma)$ with $n_{\text{eff}}$ derived, ensuring that every
configuration in the sweep adapts to distributional shifts at
the same rate---only the exploration--exploitation balance
and the granularity of forgetting differ.

We set $T_{\text{adapt}} = \hpTAdapt$ by anchoring it to the
catastrophic-failure phase length ($N_{\text{phase2}} \approx 595$
prompts in the validation split), so the router can fully override
its prior within the experimental measurement window.
In practice, a deployment engineer would set $T_{\text{adapt}}$ to the
shortest acceptable reaction time for their service (e.g., the
time to detect and reallocate away from a degraded model).

\paragraph{Multi-objective formulation.}
With $n_{\text{eff}}$ derived, we score each $(\alpha, \gamma)$
configuration on two objectives:
%
\begin{enumerate}[nosep,leftmargin=*]
  \item \textbf{Budget-Paced Pareto AUC} (stationary efficiency):
    sweep log-spaced budget targets on the validation split with the
    BudgetPacer active, build per-seed Pareto frontiers, average AUC
    across $\hpGridAlpha \times \hpGridGamma = 42$ configs $\times$ 20~seeds.

  \item \textbf{Catastrophic-failure Phase-2 reward} (non-stationary
    resilience):
    two-phase simulation on the validation split---Phase~1 (first
    half) with normal rewards, Phase~2 (second half) with the failed
    arm's reward collapsed to near-zero and cost set to \$0.
    Mean Phase-2 reward captures how quickly the router detects the
    failure and reallocates traffic.
    We tune on Mistral failure, the most operationally stressful
    scenario (Mistral is the dominant arm across the broadest range
    of budget targets), and validate generalization to all $K$ arms
    post-hoc.
\end{enumerate}

\paragraph{Pareto knee-point selection.}
Rather than scalarising the two objectives with arbitrary
weights or imposing an $\epsilon$-constraint, we build the
\emph{Pareto frontier} of non-dominated (AUC, P2\_reward)
configurations and select the \textbf{knee point}---the
configuration with maximum perpendicular distance from the
line connecting the two extreme endpoints of the frontier,
after min--max normalisation for scale invariance.
The knee point represents the inflection where improving one
objective begins to require disproportionate sacrifice of the
other, and requires no weight specification or tolerance parameter.
%% Fallbacks if bootstrap commands not yet generated:
\providecommand{\hpBootNIter}{2000}
\providecommand{\hpBootKneeFreq}{--}
\providecommand{\hpBootNeighborFreq}{--}
\providecommand{\hpBootNUnique}{--}
%%
Seed-level bootstrap resampling (\hpBootNIter{} iterations)
confirms the selection is stable: the knee-point configuration is
chosen in \hpBootKneeFreq\% of bootstrap samples, and a
configuration within one $\gamma$ grid step is chosen in
\hpBootNeighborFreq\% of samples (only \hpBootNUnique{} unique
configs are ever selected across all resamples).

\paragraph{Results.}
Table~\ref{tab:hparam_selection} summarises the selection outcome
for both ParetoBandit (warmup priors) and Tabula Rasa (cold start).

\begin{table}[t]
\centering
\caption{%
  \textbf{$T_{\text{adapt}}$-constrained Pareto knee-point selection.}
  AUC-only optimisation selects $\gamma{=}1.0$ (no forgetting)
  for both variants.
  The knee-point method selects moderate forgetting for
  \emph{both} variants: $\gamma{=}\hpParetoBanditGamma$ for
  ParetoBandit and $\gamma{=}\hpTabulaGamma$ for Tabula Rasa.
  Warmup priors give ParetoBandit higher AUC, higher P2 reward,
  and substantially lower variance than the cold-start baseline.
  Grid: $\alpha \in \{0.01, \ldots, 1.0\}$
  ($\hpGridAlpha$~values) $\times$
  $\gamma \in \{0.994, \ldots, 1.0\}$ ($\hpGridGamma$~values),
  $n_{\text{eff}}$ derived via $T_{\text{adapt}}{=}\hpTAdapt$.
  Metrics: budget-paced Pareto AUC and Phase-2 reward under
  catastrophic Mistral failure (val split, 20~seeds).%
}
\label{tab:hparam_selection}
\small
\begin{tabular}{@{}llccccc@{}}
\toprule
\textbf{Variant} & \textbf{Method}
  & $\boldsymbol{\alpha}$ & $\boldsymbol{n_{\text{eff}}}$
  & $\boldsymbol{\gamma}$
  & \textbf{BP AUC} & \textbf{P2 Reward} \\
\midrule
\multirow{2}{*}{ParetoBandit}
  & AUC-only     & $\hpAUCOnlyAlpha$ & $\hpAUCOnlyNeff$ & $\hpAUCOnlyGamma$ & $\hpAUCOnlyAUC$ & --- \\
  & Knee-point   & $\hpParetoBanditAlpha$ & $\hpParetoBanditNeff$ & $\hpParetoBanditGamma$ & $\hpParetoBanditAUC$ & $\hpParetoBanditPTwoReward$ \\[3pt]
\multirow{2}{*}{Tabula Rasa}
  & AUC-only     & $\hpTabulaAUCOnlyAlpha$ & --- & $\hpTabulaAUCOnlyGamma$ & $\hpTabulaAUCOnlyAUC$ & --- \\
  & Knee-point   & $\hpTabulaAlpha$ & --- & $\hpTabulaGamma$ & $\hpTabulaAUC$ & $\hpTabulaPTwoReward$ \\
\bottomrule
\end{tabular}
\end{table}

AUC-only optimisation selects $\gamma{=}1.0$ (no forgetting) for
both variants---confirming that single-objective tuning converges
on the stationary optimum, which
Section~\ref{sec:eval_nonstationary} shows is catastrophic under drift.
The Pareto knee-point method introduces moderate forgetting for
\emph{both} variants: $\gamma{=}\hpParetoBanditGamma$ (ParetoBandit)
and $\gamma{=}\hpTabulaGamma$ (Tabula Rasa), trading a small amount
of stationary AUC (${\sim}\hpAUCSacrifice\%$ for ParetoBandit) for
substantially improved failure resilience.

On the held-out test split, the selected configurations achieve
Pareto AUC of $\hpParetoBanditTestAUC$ (ParetoBandit) and
$\hpTabulaTestAUC$ (Tabula Rasa) versus $\hpFixedTestAUC$ for the
fixed-model envelope---deltas of $\hpParetoBanditTestDelta\%$ and
$\hpTabulaTestDelta\%$, both within one standard deviation of
seed variance ($\pm\hpParetoBanditTestStd$ and
$\pm\hpTabulaTestStd$).

\paragraph{Cross-arm validation.}
To confirm that the selected configuration generalises beyond the
tuning arm (Mistral), we evaluate it against catastrophic failure
of all $K{=}3$ arms independently.
Table~\ref{tab:cross_arm_validation} shows that Phase-2 recovery
is robust across all failure scenarios, with Mistral failure---the
tuning target---being the hardest case (lowest P2 reward), as
expected from its role as the dominant arm across the widest range
of budget targets.

\begin{table}[t]
\centering
\caption{%
  \textbf{Cross-arm failure validation} for the selected
  ParetoBandit configuration
  ($\alpha{=}\hpParetoBanditAlpha$,
   $\gamma{=}\hpParetoBanditGamma$,
   $n_{\text{eff}}{=}\hpParetoBanditNeff$).
  Phase-2 mean reward under catastrophic failure of each arm
  independently (val split, 20~seeds, 3~budget targets).
  Mistral failure (the tuning target) is the hardest scenario,
  confirming it as the appropriate arm for tuning.
  The forgetting mechanism generalises to all arms.%
}
\label{tab:cross_arm_validation}
\small
\begin{tabular}{@{}lcc@{}}
\toprule
\textbf{Failed Arm} & \textbf{P2 Reward} & \textbf{Std} \\
\midrule
Llama-8B             & $\hpCrossArmLlamaPTwo$   & $\pm\hpCrossArmLlamaStd$ \\
Mistral-Large (tuned on) & $\hpCrossArmMistralPTwo$ & $\pm\hpCrossArmMistralStd$ \\
Gemini-Pro           & $\hpCrossArmGeminiPTwo$  & $\pm\hpCrossArmGeminiStd$ \\
\bottomrule
\end{tabular}
\end{table}

Three findings merit attention:

\begin{enumerate}[nosep,leftmargin=*]
  \item \textbf{The adaptation-horizon coupling eliminates a
    hyperparameter.}
    By deriving $n_{\text{eff}}$ from $\gamma$ and $T_{\text{adapt}}$, the 3D search
    space $(\alpha, n_{\text{eff}}, \gamma)$ collapses to 2D
    $(\alpha, \gamma)$ with a single, interpretable control:
    $T_{\text{adapt}}$ specifies how many queries the router needs to
    override its prior after a shift.
    Every configuration in the sweep shares the same adaptation
    speed by construction, so the knee-point selection operates on
    a surface where all candidates can detect and respond to
    distributional changes within the same horizon.
    The practitioner tunes a single deployment-meaningful parameter
    ($T_{\text{adapt}}$) rather than two abstract statistical quantities
    ($n_{\text{eff}}$ and $\gamma$).

  \item \textbf{The knee-point method selects forgetting for both
    variants.}
    A hard-threshold approach (e.g., retaining only configurations
    within $\epsilon$ of the best AUC) would lock Tabula Rasa to
    $\gamma{=}1.0$, because even mild forgetting pushes the
    cold-start variant outside any reasonable AUC tolerance.
    The Pareto knee-point avoids this failure mode:
    it selects $\gamma{=}\hpTabulaGamma$ for Tabula Rasa by
    identifying the inflection where the marginal AUC cost of
    additional forgetting begins to accelerate, rather than
    imposing a binary feasibility cutoff.
    For Tabula Rasa, this balance point still involves some
    forgetting, even though the cold-start variant pays a steeper
    AUC price for it than ParetoBandit does.

  \item \textbf{Warmup priors remain the structural enabler.}
    While both variants now adopt forgetting, ParetoBandit's
    configuration is strictly superior on both objectives:
    higher AUC ($\hpParetoBanditAUC$ vs.\ $\hpTabulaAUC$) and
    higher P2 reward ($\hpParetoBanditPTwoReward$ vs.\
    $\hpTabulaPTwoReward$), with substantially lower variance.
    The mechanism is unchanged: under exponential discounting, the
    Discounted LinUCB estimate shrinks toward its prior.
    For Tabula Rasa, that prior is $(\lambda I, \mathbf{0})$---the
    estimate shrinks toward ignorance.
    For ParetoBandit, the prior encodes $\hpParetoBanditNeff$
    pseudo-observations per arm from the training split; forgetting
    erodes only the online residual corrections, not the baseline
    policy.
    The cross-arm validation (Table~\ref{tab:cross_arm_validation})
    further confirms that this advantage generalises: the
    forgetting mechanism is arm-agnostic, and the selected
    configuration recovers effectively regardless of which arm
    fails.
\end{enumerate}

\paragraph{Novelty and significance.}
This $T_{\text{adapt}}$-constrained Pareto knee-point protocol is, to our
knowledge, the first multi-objective hyperparameter selection
framework for non-stationary bandit routing.
It addresses a gap that spans both the LLM routing and the general
bandit literatures:
existing systems either assume stationarity in their tuning
protocol~\cite{panda2025pilot,wang2025barp,li2025llmbandit},
delegate adaptation to model retraining rather than tuning bandit
parameters~\cite{wang2025mixllm}, or derive $\gamma$ from
theoretical quantities (number of breakpoints) that are unknown in
production~\cite{garivier2011switching}.
The CDT framework~\cite{kang2024cdt} demonstrates that online
hyperparameter tuning is feasible for contextual bandits, but it
optimises a single reward objective and does not address the
stability--plasticity tradeoff that arises when the router must
simultaneously serve stationary traffic well and recover from shifts.

Our contribution is threefold:
(i)~the adaptation-horizon coupling
(Equation~\ref{eq:t_adapt}) reduces the search space from 3D to 2D
and replaces two abstract parameters ($n_{\text{eff}}$, $\gamma$) with a
single deployment-meaningful quantity ($T_{\text{adapt}}$);
(ii)~the Pareto knee-point selection is parameter-free (no weights
or tolerance to specify) and naturally identifies the balanced
trade-off between stationary efficiency and non-stationary
resilience;
(iii)~the cross-arm validation step
(Table~\ref{tab:cross_arm_validation}) provides empirical evidence
that the selected forgetting rate generalises beyond the specific
failure scenario used for tuning.

The empirical result that single-objective tuning
selects $\gamma{=}1.0$ (no forgetting) for \emph{all} variants
is itself a cautionary finding: any practitioner who tunes via
standard cross-validation will converge on no forgetting,
offering zero protection against distribution shift---despite
this being the failure mode they are most likely to encounter
in production.

\paragraph{Practical implications.}
From a deployment perspective, the protocol requires the
practitioner to specify two quantities: (i)~the adaptation horizon
$T_{\text{adapt}}$ (how quickly the router should be able to react to a
distributional shift), and (ii)~a grid of $(\alpha, \gamma)$
candidates.
The output is a single configuration that is
near-optimal for budget-paced routing and maximally adaptive within
that quality budget, with no weight or tolerance parameter to
calibrate.
Because the secondary objective uses a synthetic catastrophic-failure
simulation, the protocol requires no live non-stationary data and
can be run entirely on a static validation split before deployment.
This makes the tuning procedure safe, offline, and repeatable---a
necessary property for production systems subject to change-management
governance.
